% AUTHOR: Amlal El Mahrouss.
% PURPOSE: Draft One.

\documentclass[10pt]{article}

%% ---- Core Packages ----------------------------------------
\usepackage[T1]{fontenc}
\usepackage[utf8]{inputenc}
\usepackage{lmodern}
\usepackage{microtype}

%% ---- Page Geometry ----------------------------------------
\usepackage[
top=1in, bottom=1in,
left=0.75in, right=0.75in,
columnsep=0.25in
]{geometry}

%% ---- Mathematics ------------------------------------------
\usepackage{amsmath, amssymb, amsthm, amsfonts}
\usepackage{mathtools}
\usepackage{bm}
\usepackage{textcomp}
\usepackage{esint}

%% ---- Figures & Tables -------------------------------------
\usepackage{graphicx}
\usepackage{booktabs}
\usepackage{multirow}
\usepackage{array}
\usepackage{caption}
\usepackage{subcaption}
\usepackage{float}

%% ---- Code Listings ----------------------------------------
\usepackage{listings}
\usepackage{xcolor}
\usepackage{algorithm}
\usepackage{algpseudocode}

\lstdefinestyle{codestyle}{
	backgroundcolor=\color{gray!8},
	basicstyle=\ttfamily\footnotesize,
	breakatwhitespace=false,
	breaklines=true,
	captionpos=b,
	commentstyle=\color{green!50!black},
	keywordstyle=\color{blue}\bfseries,
	stringstyle=\color{orange!70!black},
	numberstyle=\tiny\color{gray},
	numbers=left,
	numbersep=5pt,
	showstringspaces=false,
	frame=single,
	rulecolor=\color{gray!40},
	tabsize=2
}
\lstset{style=codestyle}

%% ---- Hyperlinks -------------------------------------------
\usepackage[
colorlinks=true,
linkcolor=blue!70!black,
citecolor=green!50!black,
urlcolor=blue!60!black
]{hyperref}
\usepackage{url}

%% ---- Bibliography -----------------------------------------
\usepackage[numbers, sort&compress]{natbib}

%% ---- Miscellaneous ----------------------------------------
\usepackage{enumitem}
\usepackage{siunitx}
\usepackage{cleveref}

%% ---- Theorem Environments ---------------------------------
\theoremstyle{plain}
\newtheorem{theorem}{Theorem}[section]
\newtheorem{lemma}[theorem]{Lemma}
\newtheorem{corollary}[theorem]{Corollary}
\newtheorem{proposition}[theorem]{Proposition}

\theoremstyle{definition}
\newtheorem{definition}[theorem]{Definition}
\newtheorem{example}[theorem]{Example}

\theoremstyle{remark}
\newtheorem{remark}[theorem]{Remark}

%% ---- Custom Commands --------------------------------------
\newcommand{\R}{\mathbb{R}}
\newcommand{\N}{\mathbb{N}}
\newcommand{\Z}{\mathbb{Z}}
\newcommand{\E}{\mathbb{E}}
\newcommand{\Prob}{\mathbb{P}}
\newcommand{\norm}[1]{\left\lVert #1 \right\rVert}
\newcommand{\abs}[1]{\left\lvert #1 \right\rvert}
\newcommand{\inner}[2]{\left\langle #1,\, #2 \right\rangle}
\newcommand{\ie}{\textit{i.e.}\xspace}
\newcommand{\eg}{\textit{e.g.}\xspace}
\newcommand{\etal}{\textit{et al.}\xspace}

%% ---- Caption Formatting -----------------------------------
\captionsetup{
	font=small,
	labelfont=bf,
	format=hang,
	justification=justified
}

\renewcommand{\qedsymbol}{\ensuremath{\blacksquare}}

\usepackage{lettrine}

%% ---- Header / Footer --------------------------------------
\usepackage{fancyhdr}
\pagestyle{fancy}
\fancyhf{}
\fancyhead[R]{\small\thepage}
\renewcommand{\headrulewidth}{0.4pt}

%% ============================================================
%%  FRONT MATTER
%% ============================================================

\title{%
	\vspace{-1.5em}%
	\large\textbf{More Notes \# 1.}%
}
\author{By Amlal El Mahrouss}
\date{\today}

%% ============================================================
\begin{document}
	%% ============================================================
	
	\maketitle
	\tableofcontents
	\newpage

%% ============================================================
\section{Several Observations First:}
%% ============================================================

\begin{theorem}
There exist an equation in a domain $\mathbb{D}$ from a non-zero $x, y, t \in \mathbb{R}$.:
\begin{align*}
    \mathcal{C} &= \mathcal{N}(x) - y \cdot \int_{n}^{p} \mathcal{T}_{n} \cdot t \, \partial t,\\
    \mathcal{C} - \mathcal{N}(x) &= -y \cdot \int_{n}^{p} \mathcal{T}_{n} \cdot t \, \partial t,\\
    \mathcal{C} \pm \frac{\mathcal{N}(x)}{y} &= \int_{n}^{p} \mathcal{T}_{n} \cdot t \, \partial t, \, t \in \mathbb{R}.
\end{align*}
That yields the identity for an integral of tensors $\mathcal{T}_n$, that would converge for $p \neq \pm \infty$,\\ and yield a non-zero value.
\end{theorem}

\begin{remark}
There exist a ratio $\forall x, y \in \mathbb{R}$, $\forall p \in \mathbb{N}$:
\begin{align*}
	t_p &= \frac{\Delta x_p}{y_p!}.
\end{align*}
Consider the set of derivatives of $t_p$ below written as the following for $y_p \neq 0$ and $y_p! \neq 0$.:
\begin{align*}
	\mathcal{S} &= \{ \frac{\partial t_p}{\partial x} \cdot \frac{\Delta x_p}{y_p!}, \, \dots \}
\end{align*}
\end{remark}

\section{Several Definitions before Several Proofs:}

Going forward on this note, we shall note the following observations for our future work:

\begin{definition}
	Let us have the definition of a normal, there exists $z$, and $\forall z, x, y \in \mathbb{R}$ and that there exists $\forall i, j, k \in \mathbb{R}$.
	\begin{equation}
		\mathcal{C} + ||z|| = \frac{\partial \, x}{\partial \, y} \cdot \mathcal{X}_{ijk} \cdot x.
	\end{equation}
	There exist a normal that cancels at $z=\pi$ equating to $\frac{\partial \, x}{\partial \, y} \mathcal{X}_{ijk} \cdot x$ which when integrated is ought to be less than $||\frac{z}{x}||$. $\mathcal{C}$ is to be considered a complementary variable in $\mathbb{R}$.
\end{definition}

\begin{remark}
	$\mathcal{X}_{ijk}$ shall be a non-zero variable factoring $x$.
\end{remark}

\begin{remark}
\label{rem:math-special-case}
	One may separately note the following for a analytical two dimensional space $\mathcal{D}$ for the function $\mathcal{T}$ and tensors $\mathcal{X}$:
	\begin{equation}
		\mathcal{T}_0 + \mathcal{R}_{t} = \iint_{\mathcal{D}} \mathcal{T}(\mathcal{X}_{x}, \mathcal{X}_{y})_{t} \, dx \land dy.
	\end{equation}
	Then:
	\begin{equation}
		\mathcal{T}_0 = \iint_{\mathcal{D}} \mathcal{T}(\mathcal{X}_{x}, \mathcal{X}_{y})_{t} \, dx \land dy - \mathcal{R}_{t}.
	\end{equation}
	A general use case for a strictly positive natural $n$ and simulation variable $x$ is:
	\begin{equation}
		\mathcal{T}_{0+n} = \iint_{\mathcal{D}} \mathcal{T}(\mathcal{X}_{x+n}, \mathcal{X}_{y+n})_{t+x} \, dx \land dy - \mathcal{R}_{t+n}.
	\end{equation}
	Or, more elegantly:
	\begin{equation}
		\mathcal{T}_{n} = \iint_{\mathcal{D}} \mathcal{T}(\mathcal{X}_{x+n}, \mathcal{X}_{y+n})_{t+x} \, dx \land dy - \mathcal{R}_{t+n}.
	\end{equation}
	Where $\forall x, y \in \mathbb{R}$.
\end{remark}

\begin{corollary}
	$\mathcal{T}_{n}'$ is defined as per $\operatorname{Trans}$ being the transform operator of $t|x$:
	\begin{equation}
		\mathcal{T}_{n}' =\operatorname{Trans}_{t|x}(\iint_{\mathcal{D}} \mathcal{T}(\mathcal{X}_{x+n}, \mathcal{X}_{y+n})_{t+x} \, dx \land dy - \mathcal{R}_{t+n}).
	\end{equation}
	Per Remark \ref{rem:math-special-case} where $t|x$ is to be read as "time $t$ over simulation variable $x$".
\end{corollary}

\nocite{*}
\bibliographystyle{acm}
\bibliography{weil1967basic, pinter2010book, S0022314X10001836, Fatou1906, A_Course_of_Pure_Mathematics, General_Theory_Dirichlet_Series, Riemann, bonito2026, hardy1922theory, chambadal1981dictionnaire}

%% ============================================================
\end{document}
%% ============================================================